\documentclass[11pt]{article}
\input{style-sujet}
\usepackage{array}

\newcommand{\note}[1]{\par\smallskip{\itshape\small #1}\par\smallskip}

\begin{document}

\entetesujet{Sujet bac 2016 -- Série C}

% ====================== EXERCICE 1 ======================
\exo{1}{4 points}

On donne dans $\Z$ l'équation :
\[ (E) : 2688x + 3024y = -3360 \]

\begin{enumerate}
  \item Déterminer le PGCD$(2688, 3024)$, puis en déduire que l'équation $(E)$ admet des solutions dans $\Z^2$.
  \item Montrer que l'équation $(E)$ est équivalente à l'équation $(E_1) : 8x + 9y = -10$.
  \item \begin{enumerate}[label=\alph*]
    \item Montrer que l'équation $(E_1)$ peut s'écrire $(E_2) : 8x \equiv -10\ [9]$.
    \item Résoudre l'équation $(E_2)$.
    \item En déduire les solutions de l'équation $(E)$.
  \end{enumerate}
\end{enumerate}

% ====================== EXERCICE 2 ======================
\exo{2}{8 points}

Le plan est orienté. $ABCD$ est un rectangle tel que $(\vv{BD},\vv{BA}) = \dfrac{\pi}{3}\ [2\pi]$. On considère le losange $BDEG$ tel que $(\vv{BD},\vv{BA}) = (\vv{BD},\vv{BE})\ [2\pi]$.

Dans cet exercice, $S_\Delta$ et $t_{\vec{u}}$ désignent respectivement la symétrie orthogonale d'axe la droite $(\Delta)$ et la translation de vecteur $\vec{u}$. On considère la transformation
\[ f = S_{(AD)} \circ S_{(AB)} \circ S_{(BD)} \]

\begin{enumerate}
  \item Faire la figure. On prendra $AB = 4$ cm et on disposera $(AB)$ horizontalement.
  \item Déterminer la nature et les éléments caractéristiques de la transformation $g = S_{(AB)} \circ S_{(BD)}$.
  \item $R$ désigne la rotation de centre $B$ et d'angle de mesure $\dfrac{\pi}{3}$. Déterminer la nature exacte de la transformation $S_{(AD)} \circ R$.
  \item On désigne par $F$ le milieu du segment $[BG]$.
  \begin{enumerate}[label=\alph*.]
    \item Démontrer que $f = t_{\vv{BF}} \circ S_{(AC)}$.
    \note{Erreur dans l'énoncé : il s'agit plutôt de montrer $S_{(AD)} \circ R = t_{\vv{BF}} \circ S_{(AC)}$. En effet, la transformation $f$ n'est pas égale à $t_{\vv{BF}} \circ S_{(AC)}$.}
    \item En déduire les éléments caractéristiques de $f$.
    \note{Substituer cette question par : Déduire les éléments caractéristiques de $S_{(AD)} \circ R$.}
  \end{enumerate}
  \item Soit $(P)$ la parabole dont une tangente est la droite $(EF)$, la normale associée est la droite $(GB)$ et l'axe focal est la droite $(EB)$. Démontrer que $A$ est le foyer de la parabole $(P)$.
  \item Soit $H$ le milieu du segment $[EG]$ et $L$ celui du segment $[ED]$. Déterminer la directrice $(d)$ de la parabole $(P)$.
  \item Construire le point $I$ de $(P)$ tel que $[IF]$ soit une corde focale.
  \item Construire la corde focale $[JK]$ de $(P)$ où $J$ appartient au segment $[AD]$.
  \item Construire l'arc d'extrémités $J$ et $F$ de $(P)$.
  \item Soit $(P')$ l'image de $(P)$ par la transformation $f$.
  \begin{enumerate}[label=\alph*.]
    \item Déterminer le foyer de la parabole $(P')$.
    \item Déterminer l'axe focal de $(P')$.
  \end{enumerate}
\end{enumerate}

% ====================== EXERCICE 3 ======================
\exo{3}{5 points}

\begin{enumerate}
  \item Déterminer la solution particulière $f$ de l'équation différentielle :
  \[ (E) : y'' + 2y' + 2y = 0 \]
  vérifiant les conditions initiales suivantes : $f\!\left(\dfrac{\pi}{2}\right) = e^{-\frac{\pi}{2}}$ et $f'\!\left(\dfrac{\pi}{2}\right) = -e^{-\frac{\pi}{2}}$.
  \item On pose $f(x) = e^{-x}\sin x$.
  \begin{enumerate}[label=\alph*.]
    \item Déterminer les réels $A$ et $B$ pour que $F(x) = e^{-x}(A\cos x + B\sin x)$ soit une primitive de $f$ sur $\R$.
    \item Calculer l'intégrale $\displaystyle\int_{n\pi}^{(n+1)\pi} f(x)\,dx$.
  \end{enumerate}
  \item Soit $(v_n)$ la suite numérique définie pour tout entier naturel $n$ par :
  \[ v_n = \frac{e^{-\pi} + 1}{2}(-e^{-\pi})^n\ ;\quad n \in \N \]
  \begin{enumerate}[label=\alph*.]
    \item Montrer que $(v_n)$ est une suite géométrique dont on précisera la raison et le premier terme.
    \item Calculer la somme des termes $S_n = v_0 + v_1 + \cdots + v_n$ en fonction de $n$, puis en déduire la limite de $S_n$ en $+\infty$.
  \end{enumerate}
\end{enumerate}

% ====================== EXERCICE 4 ======================
\exo{4}{3 points}

Soit le tableau statistique à double entrée suivant :

\begin{center}
\renewcommand{\arraystretch}{1.3}
\begin{tabular}{|c|c|c|c|}
\hline
$X\,\backslash\,Y$ & -1 & 2 & 3 \\
\hline
1 & 2 & 1 & 1 \\
\hline
2 & 2 & 3 & 1 \\
\hline
\end{tabular}
\end{center}

\begin{enumerate}
  \item Convertir ce tableau en un tableau linéaire.
  \item Déterminer le coefficient de corrélation $\rho_{X,Y}$ des caractères $X$ et $Y$. On donne $\overline{X} = 1{,}6$ et $\overline{Y} = 1$.
  \item Donner une interprétation géométrique de cette corrélation.
\end{enumerate}

\end{document}
